\section{Limits of the result and evidence gaps}\label{sec:limits}
\begin{eibox}{What this result is not}
It is a computational design (\gStatus). Nothing has been printed, stress-relieved, ground or physically tested, and no
safety factor reported here is a qualification for use.
\end{eibox}
\begin{itemize}
\item \textbf{Outcome against the targets.} The best individual \gDesign{} has worst safety-factor margin ratio
$\mu = \gWorstMarginRatio$ (on \gWorstMarginCase, multiplier $W = \gW$), finished partial-volume mass \gMassFinishedG{}~g
(\gMassG{}~g on the $0.20$~mm mesh, which \gMassReading), and sealed internal void volume \gSealedVoid{}~$\mathrm{mm^3}$
(printed solid in the build field, $d_\mathrm{am} = -h$, so the mass counts solid titanium rather than trapped powder).
\item \textbf{Blade surface texture (frozen \code{G5.0} decoder limit).} The frozen \code{G5.0} decoder allows the evolved
surface texture to reach the blade zone (for \gDesign, texture depth \code{tex\_amp} $= \gTexAmp$~mm, pitch $\gTexPitch$~mm, and
\code{scope\_blade} $= \gScopeBladeThree$), carving inward across both the blade flats and the bevels so the as-decoded edge is
notched. The edge cross-sections (Fig.~\ref{fig:edges}) and edge measurements evaluate the exact bevel law, which post-build
grinding restores because the \code{RS-karambit-v5} finishing allowance ($\gReqEdgeFinishMm$~mm per ground face) exceeds the
$\gTexAmp$~mm texture depth. Restricting surface texture outside the bevels is deferred to a later decoder generation.
\item \textbf{One platform of record.} v5.0.0's \code{arm64} records agreed closely with, but were not bit-identical to,
\code{x86\_64} re-evaluations, because the platforms' maths libraries differ in the last bits. v\gRecord{} runs the search, every
re-evaluation and the ranking on one platform (static \code{x86\_64} binaries, natively or in a \code{linux/amd64} container): every archive
row re-scores bit for bit, every pool is recomputed from its sources, \gReproIsland{} re-run on \gReproHost{} reproduced all
\gReproEntries{} logged history entries over \gReproEvals{} evaluations, and re-evaluating all \gCohortElites{} elites of the final
pool reproduced all \gTierOneReproduced{} Tier-1 rows. A different libm (another platform) would again agree closely rather than
exactly.
\item \textbf{The search grid can rank designs it cannot keep.} The search scores at $h = \gH$~mm with a section model of the tip;
the ranking re-scores the shortlist at $h = \gHRank$~mm with the solid tip and the winner must also pass the grid check and every
load case on every plate mesh. A design that leads the search grid can therefore rank lower, which is why the search-grid best of the
lineage (\gPoolFinalBest) is not the score of record (\gScore).
\item \textbf{Analytic hand-zone models, checked by FE.} The score models grip squeeze, guard push and the serration teeth
analytically. Against an independent hex8 voxel FE of the same design (Table~\ref{tab:hex8}), the score's safety factor is
\gFemGripCmp{} the FE for grip squeeze (\gFemModelGrip{} against \gFemHexGrip), \gFemGuardCmp{} it for guard push (\gFemModelGuard{}
against \gFemHexGuard) and \gFemTeethCmp{} it for the teeth (\gFemModelTeeth{} against \gFemHexTeeth); a model above the FE is
optimistic. The lowest FE safety factor, \gFemHexMin, \gFemHexClears{} the $\gReqSfYield$ requirement. The FE checks are reported,
not scored, so a later decoder could score these cases on the solid.
\item \textbf{Linear-elastic surrogates and solid-tip readout windows.} All structural evaluations are linear-elastic. Permanent set
under the GOST~R~51501-99 deflection test~\cite{gost51501} is represented by no yield at $\mathrm{SF} = \gReqSfYield$, not by an
elastoplastic solve. The 3D \code{TET10} tip submodel reports $\mathrm{SF} = \sigma_\mathrm{allow} / \sigma_{\mathrm{vM},\max}$
from the nodal maximum von Mises stress within each case's readout window ($r \in [\gTipAxFromMm, 10.0]$~mm for axial,
$[\gTipAxFromMm, \gTipLatHiMm]$~mm for side load, and $[\gTipTwReadLoMm, 15.0]$~mm for twist) to stay clear of the loaded
patch ($0.75$~mm axial/lateral, $1.0$~mm twist) and the clamped cut plane ($r_\mathrm{cut} = 30$~mm).
\item \textbf{Homogenised TPMS core and first-order hand patch.} In the blade and grip section models, the sheet-TPMS lattice
is represented by the Naghavi \emph{et al.}~\cite{naghavi2022} and Timercan \emph{et al.}~\cite{timercan2021} compression
scaling laws; periodic unit-cell finite-element homogenisation and tension-fatigue testing of the printed sheet lattice are not
yet performed. Hand contact pressure uses a first-order patch area ($60\,\%$ of grip length times half the swell perimeter)
rather than a deformable soft-tissue contact solve.
\item \textbf{Evidence gaps requiring physical coupon and part tests (\code{PROVENANCE-v5.md}).} As documented in
\src{requirements/PROVENANCE-v5.md}{\code{PROVENANCE-v5.md}}, ten domain inputs rest on derived bounds or single-orientation
studies and require physical test replacement before any fitness-for-use claim: per-tooth snag/pry loads ($\gReqSnagN/\gReqPryN$~N),
off-axis tip side/twist shares ($\gReqTipLatN$~N, $\gReqTipTwNm$~N$\cdot$m), ISO~8442-5~\cite{iso8442-5} edge retention of ground L-PBF
Ti-6Al-4V, little-finger ring avulsion hazard~\cite{hidh}, wet skin-on-titanium friction ($\mu = 0.15$), down-skin thin-wall
roughness depth ($T_s = 0.10$~mm), $80\,^\circ\mathrm{C}$ L-PBF tensile derating, and the $1.22$~m MIL-STD-810H transit drop
test~\cite{mil810h,fl1}.
\end{itemize}
